% GENERATED FILE. Edit canonical lessons, then run npm run generate:books.
% canonical-source-hash: fnv1a64:eca65e607696e524
% canonical-lessons: TE-C489-jebudonga, TE-C489-kalti, TE-C489-mosapovu, TE-C489-kullu, TE-C489-migulu

\chapter{Cheated at the Market}
\label{ch:te-cheated-at-the-market}
\hlchaptermodality{489}

\begin{chapteropening}
\textbf{By the end of this chapter:} \emph{I can say a pickpocket, adulteration, to be cheated, to rot and to be left over.}

Complete the last lesson of chapter 489: I can say a pickpocket, adulteration, to be cheated, to rot and to be left over.
\end{chapteropening}

\section[\texorpdfstring{\te{జేబుదొంగ} (jēbudoṅga)}{జేబుదొంగ (jēbudoṅga)}]{\te{జేబుదొంగ} (jēbudoṅga) — a pickpocket}

\label{lesson:TE-C489-jebudonga}



\begin{quote}
Before the new one: say the Telugu for a shortage, then the Telugu for an auction.
\end{quote}

\subsection*{You'll want to know: \te{జేబుదొంగ}}
\textbf{\te{జేబుదొంగ}} — \emph{jēbudoṅga} — \textquotedblleft{}a pickpocket\textquotedblright{}.

\begin{quote}
\te{రద్దీ} \te{బస్సులో} \te{జేబుదొంగ} \te{నా} \te{పర్సు} \te{కొట్టేశాడు}.
\end{quote}

\begin{quote}
\emph{raddī bassulō jēbudoṅga nā parsu koṭṭēśāḍu.}
\end{quote}

\begin{quote}
In the crowded bus a pickpocket took my purse.
\end{quote}

\te{జేబు}, a pocket, + \te{దొంగ}, a thief.

\begin{cousinweb}
\textbf{\te{జేబుదొంగ}} (\emph{jēbudoṅga}) is a Telugu compound: \te{జేబు}, \textquotedblleft{}a pocket\textquotedblright{}, a loan from Arabic \emph{jayb} through Persian and Hindustani, and \te{దొంగ}, \textquotedblleft{}a thief\textquotedblright{}, native Dravidian (DEDR 2982).
\end{cousinweb}

\begin{grammarlens}[title={What you've built}]
One more word for getting around.
\end{grammarlens}

\subsection*{Your turn}
\begin{itemize}
\raggedright
  \item \emph{Say it:} \emph{jēbudoṅga}
  \item \emph{Say it:} \emph{jēbudoṅga}, once more
  \item \emph{Say it:} \emph{raddī bassulō jēbudoṅga nā parsu koṭṭēśāḍu.} — In the crowded bus a pickpocket took my purse.
  \item \emph{Recall:} say the Telugu for a shortage, then the Telugu for an auction, then say \emph{jēbudoṅga} again
\end{itemize}

\begin{tcolorbox}[breakable,colback=teal!4,colframe=teal!35!black,title={Before you move on}]
What does \te{జేబుదొంగ} mean? (\textquotedblleft{}A pickpocket\textquotedblright{}.) Say it once more.
\end{tcolorbox}
\section[\texorpdfstring{\te{కల్తీ} (kaltī)}{కల్తీ (kaltī)}]{\te{కల్తీ} (kaltī) — adulteration}

\label{lesson:TE-C489-kalti}



\begin{quote}
Before the new one: say the Telugu for an auction, then the Telugu for a pickpocket.
\end{quote}

\subsection*{You'll want to know: \te{కల్తీ}}
\textbf{\te{కల్తీ}} — \emph{kaltī} — \textquotedblleft{}adulteration\textquotedblright{}.

\begin{quote}
\te{ఆ} \te{పాలలో} \te{కల్తీ} \te{ఉందని} \te{తేలింది}.
\end{quote}

\begin{quote}
\emph{ā pālalō kaltī undani tēlindi.}
\end{quote}

\begin{quote}
It turned out that the milk was adulterated.
\end{quote}

\te{కల్తీ} \te{చేయు}, to adulterate; \te{కల్తీ} \te{నూనె}, adulterated oil.

\begin{cousinweb}
\textbf{\te{కల్తీ}} (\emph{kaltī}) is a Telugu word for adulteration whose origin is not settled.
\end{cousinweb}

\begin{grammarlens}[title={What you've built}]
One more word for getting around.
\end{grammarlens}

\subsection*{Your turn}
\begin{itemize}
\raggedright
  \item \emph{Say it:} \emph{kaltī}
  \item \emph{Say it:} \emph{kaltī}, once more
  \item \emph{Say it:} \emph{ā pālalō kaltī undani tēlindi.} — It turned out that the milk was adulterated.
  \item \emph{Recall:} say the Telugu for an auction, then the Telugu for a pickpocket, then say \emph{kaltī} again
\end{itemize}

\begin{tcolorbox}[breakable,colback=teal!4,colframe=teal!35!black,title={Before you move on}]
What does \te{కల్తీ} mean? (\textquotedblleft{}Adulteration\textquotedblright{}.) Say it once more.
\end{tcolorbox}
\section[\texorpdfstring{\te{మోసపోవు} (mōsapōvu)}{మోసపోవు (mōsapōvu)}]{\te{మోసపోవు} (mōsapōvu) — to be cheated}

\label{lesson:TE-C489-mosapovu}



\begin{quote}
Before the new one: say the Telugu for a pickpocket, then the Telugu for adulteration.
\end{quote}

\subsection*{You'll want to know: \te{మోసపోవు}}
\textbf{\te{మోసపోవు}} — \emph{mōsapōvu} — \textquotedblleft{}to be cheated\textquotedblright{}.

\begin{quote}
\te{నకిలీ} \te{బంగారం} \te{కొని} \te{అతను} \te{మోసపోయాడు}.
\end{quote}

\begin{quote}
\emph{nakilī baṅgāraṁ koni atanu mōsapōyāḍu.}
\end{quote}

\begin{quote}
He was cheated when he bought fake gold.
\end{quote}

Past (I …): \te{మోసపోయాను}, \emph{mōsapōyānu}.

\te{మోసం}, a fraud, + \te{పోవు}: to fall into a fraud. \te{మోసం} \te{చేయు} is to cheat someone.

\begin{cousinweb}
\textbf{\te{మోసపోవు}} (\emph{mōsapōvu}) is a Telugu compound: \te{మోసం}, \textquotedblleft{}deceit\textquotedblright{}, usually traced to Sanskrit \emph{mōṣa}, \textquotedblleft{}theft\textquotedblright{} (CDIAL 10358), and the native \te{పోవు}, \textquotedblleft{}to go\textquotedblright{} (DEDR 4572).
\end{cousinweb}

\begin{grammarlens}[title={What you've built}]
One more word for getting around.
\end{grammarlens}

\subsection*{Your turn}
\begin{itemize}
\raggedright
  \item \emph{Say it:} \emph{mōsapōvu}
  \item \emph{Say it:} \emph{mōsapōvu}, once more
  \item \emph{Say it:} \emph{nakilī baṅgāraṁ koni atanu mōsapōyāḍu.} — He was cheated when he bought fake gold.
  \item \emph{Recall:} say the Telugu for a pickpocket, then the Telugu for adulteration, then say \emph{mōsapōvu} again
\end{itemize}

\begin{tcolorbox}[breakable,colback=teal!4,colframe=teal!35!black,title={Before you move on}]
What does \te{మోసపోవు} mean? (\textquotedblleft{}To be cheated\textquotedblright{}.) Say it once more.
\end{tcolorbox}
\section[\texorpdfstring{\te{కుళ్ళు} (kuḷḷu)}{కుళ్ళు (kuḷḷu)}]{\te{కుళ్ళు} (kuḷḷu) — to rot}

\label{lesson:TE-C489-kullu}



\begin{quote}
Before the new one: say the Telugu for adulteration, then the Telugu for to be cheated.
\end{quote}

\subsection*{You'll want to know: \te{కుళ్ళు}}
\textbf{\te{కుళ్ళు}} — \emph{kuḷḷu} — \textquotedblleft{}to rot\textquotedblright{}.

\begin{quote}
\te{వేడికి} \te{పండ్లు} \te{రెండు} \te{రోజుల్లో} \te{కుళ్ళాయి}.
\end{quote}

\begin{quote}
\emph{vēḍiki paṇḍlu reṇḍu rōjullō kuḷḷāyi.}
\end{quote}

\begin{quote}
In the heat the fruit rotted in two days.
\end{quote}

Past (it …): \te{కుళ్ళింది}, \emph{kuḷḷindi}.

\te{కుళ్ళిన}, rotten: \te{కుళ్ళిన} \te{టమాటాలు}, rotten tomatoes. Its subject is a thing, so the past shown is the third person: \te{కుళ్ళింది}, it rotted.

\begin{cousinweb}
\textbf{\te{కుళ్ళు}} (\emph{kuḷḷu}) is a native-looking Telugu verb whose origin is not settled.
\end{cousinweb}

\begin{grammarlens}[title={What you've built}]
One more word for getting around.
\end{grammarlens}

\subsection*{Your turn}
\begin{itemize}
\raggedright
  \item \emph{Say it:} \emph{kuḷḷu}
  \item \emph{Say it:} \emph{kuḷḷu}, once more
  \item \emph{Say it:} \emph{vēḍiki paṇḍlu reṇḍu rōjullō kuḷḷāyi.} — In the heat the fruit rotted in two days.
  \item \emph{Recall:} say the Telugu for adulteration, then the Telugu for to be cheated, then say \emph{kuḷḷu} again
\end{itemize}

\begin{tcolorbox}[breakable,colback=teal!4,colframe=teal!35!black,title={Before you move on}]
What does \te{కుళ్ళు} mean? (\textquotedblleft{}To rot\textquotedblright{}.) Say it once more.
\end{tcolorbox}
\section[\texorpdfstring{\te{మిగులు} (migulu)}{మిగులు (migulu)}]{\te{మిగులు} (migulu) — to be left over}

\label{lesson:TE-C489-migulu}



\begin{quote}
Before the new one: say the Telugu for to be cheated, then the Telugu for to rot.
\end{quote}

\subsection*{You'll want to know: \te{మిగులు}}
\textbf{\te{మిగులు}} — \emph{migulu} — \textquotedblleft{}to be left over\textquotedblright{}.

\begin{quote}
\te{పెళ్ళి} \te{తరువాత} \te{చాలా} \te{భోజనం} \te{మిగిలింది}.
\end{quote}

\begin{quote}
\emph{peḷḷi taruvāta cālā bhōjanaṁ migilindi.}
\end{quote}

\begin{quote}
After the wedding a lot of food was left over.
\end{quote}

Past (I …): \te{మిగిలాను}, \emph{migilānu}.

What is left takes no ending; the one who keeps it takes -\te{కి}: \te{నాకు} \te{వంద} \te{రూపాయలు} \te{మిగిలాయి}.

\begin{cousinweb}
\textbf{\te{మిగులు}} (\emph{migulu}) is native Dravidian: DEDR 4838 lists Telugu \emph{migulu}, \textquotedblleft{}to remain, to be left over, to be saved\textquotedblright{}, with Tamil \emph{miku}, \textquotedblleft{}to exceed\textquotedblright{}.
\end{cousinweb}

\begin{grammarlens}[title={What you've built}]
One more word for getting around.
\end{grammarlens}

\subsection*{Your turn}
\begin{itemize}
\raggedright
  \item \emph{Say it:} \emph{migulu}
  \item \emph{Say it:} \emph{migulu}, once more
  \item \emph{Say it:} \emph{peḷḷi taruvāta cālā bhōjanaṁ migilindi.} — After the wedding a lot of food was left over.
  \item \emph{Recall:} say the Telugu for to be cheated, then the Telugu for to rot, then say \emph{migulu} again
\end{itemize}

\begin{tcolorbox}[breakable,colback=teal!4,colframe=teal!35!black,title={Before you move on}]
What does \te{మిగులు} mean? (\textquotedblleft{}To be left over\textquotedblright{}.) Say it once more.
\end{tcolorbox}
